
Le @@M0@@-anneau @@M1@@ est caractérisé par la propriété universelle suivante :

\textbf{Proposition 4.2.} \uline{Pour toute} @@M2@@-@@M3@@-\uline{algèbre} @@M4@@, \uline{et toute famille} @@M5@@ \uline{d'éléments de} @@M6@@, \uline{il existe un} @@M7@@-@@M8@@-\uline{homomorphisme et un seul} @@M9@@ \uline{de} @@M10@@ \uline{dans} @@M11@@ \uline{tel que pour tout} @@M12@@, @@M13@@.

Pour que @@M14@@ soit un @@M15@@-homomorphisme, il faut que @@M16@@ pour tous @@M17@@~; ceci définit @@M18@@ de façon unique. L'homomorphisme ainsi obtenu est bien un @@M19@@-homomorphisme, car pour vérifier que @@M20@@ il suffit de le vérifier sur un système de générateurs de @@M21@@ sur @@M22@@, ce qui résulte de la définition de @@M23@@ et des relations (2.4.2).

4.3. On peut faire la même construction en utilisant les @@M24@@ au lieu des @@M25@@. Soit @@M26@@~; pour définir la structure de @@M27@@-anneau de @@M28@@, on se donne un homomorphisme d'anneaux : @@M29@@, par :
@@M30@@
On vérifie ensuite la commutativité de (3.7.2), ce qui montre qu'on obtient bien une structure de @@M31@@-anneau.

Le @@M32@@-anneau @@M33@@ vérifie également la propriété universelle 4.2~; il est donc canoniquement isomorphe (en tant que @@M34@@-anneau) au @@M35@@-anneau @@M36@@.

\textbf{Définition 4.4.} \uline{Le} @@M37@@-\uline{anneau} @@M38@@ \uline{sera appelé le} @@M39@@-anneau libre engendré sur @@M40@@ par la famille de générateurs @@M41@@.

\textbf{Définition 4.5.} \uline{Soient} @@M42@@ \uline{un} @@M43@@-\uline{anneau}, @@M44@@ \uline{un idéal de} @@M45@@. \uline{On dit que} @@M46@@ \uline{est} \uline{un} @@M47@@-idéal de @@M48@@ \uline{si pour tout} @@M49@@, \uline{et tout} @@M50@@, @@M51@@.

Si @@M52@@ est un @@M53@@-homomorphisme, le noyau de @@M54@@ est un @@M55@@-idéal.

Soit @@M56@@ un @@M57@@-idéal~; l'homomorphisme composé : @@M58@@ s'annule sur @@M59@@ et se factorise donc par @@M60@@, qui se trouve ainsi muni 
